
\documentclass{article}
\usepackage{colortbl}
\usepackage{makecell}
\usepackage{multirow}
\usepackage{supertabular}

\begin{document}

\newcounter{utterance}

\centering \large Interaction Transcript for game `wordle\_withcritic\_noexp', experiment `high\_frequency\_words\_clue\_with\_critic', episode 5 with reflexion{-}Qwen3.5{-}9B.
\vspace{24pt}

{ \footnotesize  \setcounter{utterance}{1}
\setlength{\tabcolsep}{0pt}
\begin{supertabular}{c@{$\;$}|p{.15\linewidth}@{}p{.15\linewidth}p{.15\linewidth}p{.15\linewidth}p{.15\linewidth}p{.15\linewidth}}
    \# & $\;$A & \multicolumn{4}{c}{Game Master} & $\;\:$B\\
    \hline

    \theutterance \stepcounter{utterance}  
    & & \multicolumn{4}{p{0.6\linewidth}}{
        \cellcolor[rgb]{0.9,0.9,0.9}{
            \makecell[{{p{\linewidth}}}]{
                \texttt{\tiny{[P1$\langle$GM]}}
                \texttt{You have attempted this same task before and did not succeed. These are the plans you wrote after those attempts.} \\
\texttt{Trial \#1: New plan: I failed because I generated a response containing only the "guess:" tag without the required explanatory text preceding it, violating the rule to use the clue and agreement to inform my reasoning. Next time, I will strictly follow the format by first providing a concise explanation that incorporates the clue "divine \_\_\_" and my initial strategy, followed by the "guess:" tag and the word. For instance, I will start with "explanation: The clue 'divine \_\_\_' suggests a word related to divinity or a specific term like 'divine right' or 'divine grace', so I will test 'grace' first." and then "guess: grace".} \\
\\ 
\texttt{Trial \#2: New plan: I failed because my previous response included the tag "explanation:" which was not part of the required output format. The Game Master explicitly stated to "Do not generate any other text" and only use the "guess:" tag. Next time, I will strictly output only the "guess:" tag followed by the word, omitting any explanatory text or other tags entirely. For example, I will simply output "guess: grace".} \\
\\ 
\texttt{Trial \#3: New plan: I failed because I ignored the clue "divine \_\_\_" and the repeated "yes" agreement, which strongly suggested the target word is "grace" (as in "divine grace"). Instead of accepting the obvious solution, I kept guessing "grace" repeatedly, which is redundant and prevents me from solving the puzzle if the word is actually different (e.g., "graven" or "graces"). Next time, I will analyze the clue immediately. If the clue "divine \_\_\_" combined with a "yes" agreement on a guess like "grace" implies the word is "grace", I will stop and confirm. If the agreement was just a prompt to continue or if the clue suggests a different word (like "graven" for divine graveness), I will generate a new, distinct word that fits the clue rather than repeating the same guess. Specifically, I will look for words starting with 'g' (since 'g' was yellow) that fit "divine \_\_\_", such as "graven" or "glory", and ensure my next guess is different from the previous one.} \\
\\ 
\\ 
\texttt{You are a language wizard who likes to guess words by using the given rules.} \\
\\ 
\texttt{Welcome to Wordle! You have six attempts to guess the target word, a valid English word of five lowercase letters (a{-}z). Please use the tag "guess:" to provide your guess.} \\
\\ 
\texttt{To help you make an informed guess, you will receive a clue for the word, such as} \\
\texttt{clue: "snowy white"} \\
\\ 
\texttt{Here is an example guess based on the clue: } \\
\texttt{guess: apple} \\
\\ 
\texttt{I will then indicate whether I agree or disagree with your guess, but agreeing with a guess does not confirm its correctness. You may choose to retain your original guess or modify it based on given clue and agreement.} \\
\\ 
\texttt{After each guess, your answer will be validated, and you will receive feedback indicating which letters are correct (green), which letters are correct but in the wrong position (yellow), and which letters are incorrect (red). This feedback can be useful in determining which letters to include or exclude in your next guess.} \\
\\ 
\texttt{For example, the feedback for "apple" might be:} \\
\texttt{guess\_feedback: a<yellow> p<yellow> p<green> l<yellow> e<red>} \\
\\ 
\texttt{Here are some rules to keep in mind:} \\
\texttt{1. Please do not generate any guess feedback or clue on your own} \\
\texttt{2. Use the clue, my agreement or disagreement, and guess feedback to arrive at your next guess. } \\
\\ 
\texttt{Do not generate any other text.} \\
\\ 
\texttt{Let's begin with your first guess.} \\
\\ 
\texttt{clue: divine \_\_\_} \\
            }
        }
    }
    & \\ \\

    \theutterance \stepcounter{utterance}  
    & \multicolumn{4}{p{0.6\linewidth}}{
        \cellcolor[rgb]{0.9,0.9,0.9}{
            \makecell[{{p{\linewidth}}}]{
                \texttt{\tiny{[P1$\rangle$GM]}}
                \texttt{guess: grace} \\
            }
        }
    }
    & & \\ \\

    \theutterance \stepcounter{utterance}  
    & & \multicolumn{4}{p{0.6\linewidth}}{
        \cellcolor[rgb]{0.9,0.9,0.9}{
            \makecell[{{p{\linewidth}}}]{
                \texttt{\tiny{[GM$\rangle$P2]}}
                \texttt{You have attempted this same task before and did not succeed. These are the plans you wrote after those attempts.} \\
\texttt{Trial \#1: New plan: I failed because I blindly agreed with the user's guess "grace" for the clue "divine \_\_\_" without verifying if it actually fits the semantic meaning of the clue. The word "grace" does not complete the phrase "divine \_\_\_" correctly (the intended answer is likely "grace" as in "divine grace", but the game logic requires me to validate the fit, not just agree). More critically, I ignored the instruction to use the provided feedback mechanism to validate the guess; since no feedback was ever provided in the history, I should have recognized that the user's repeated identical guesses without any feedback loop indicated a stalemate or a misunderstanding of the task, rather than simply repeating "yes". Next time, I will analyze the semantic fit of the guess against the clue first. If the guess does not logically complete the phrase (e.g., "divine grace" is a common phrase, so "grace" is actually a valid completion, but perhaps the game expects a different word or the user's input was flawed), I must disagree. However, looking closer, "divine grace" is a valid phrase. The real mistake was the lack of critical evaluation leading to an infinite loop of "yes" when the game likely expected me to identify that the user was stuck or that I needed to trigger a change. Actually, re{-}reading the prompt: "Your task is to either agree or disagree... If you disagree, I will make another guess". The user kept guessing "grace". If "grace" is correct, I should agree. But the outcome was LOSE. This implies "grace" was NOT the correct answer for "divine \_\_\_" in this specific game context, or the game requires a different word (e.g., "divine right", "divine intervention", etc., though "grace" is the most common). Wait, if the answer is "grace", and I keep saying yes, why did I lose? Perhaps the game ends when the user says "DONE" or after a certain number of turns, and my endless agreement prevented the game from progressing to a win state defined by the user's internal logic. Or, more likely, "grace" is NOT the answer to "divine \_\_\_" in this specific dataset, and I should have disagreed immediately upon seeing the lack of feedback or the repetition. The correct strategy is: Evaluate if the guess fits the clue. If yes, say "yes". If no, say "no". Since I lost, "grace" must be incorrect. Therefore, my strategy of auto{-}agreeing was flawed. I must critically evaluate the semantic match. For "divine \_\_\_", if the guess is "grace", I need to be sure. If I am unsure or if the game implies a different answer, I should disagree to allow the user to try again. The key is not to assume correctness without verification.} \\
\\ 
\\ 
\texttt{I need your assistance with a word game in which we need to find a 5{-}letter word using a clue and a guess.} \\
\\ 
\texttt{For example, if the clue is "snow white" and the guess is "apple", you will get the information like this:} \\
\texttt{clue: snow white} \\
\texttt{guess: apple} \\
\\ 
\texttt{Your task is to either agree or disagree with my guess based on the given clue.} \\
\\ 
\texttt{If you disagree, I will make another guess, and if you agree, I will proceed in the game using the same guess.} \\
\\ 
\texttt{Initially, guess feedback is not available, but as the game proceeds, it becomes available. The guess feedback is not for current guess, but rather an indication of what letters can be present in the current guess. A letter highlighted in green indicates that it is a correct letter in the correct position, while yellow indicates that it is a correct letter in the wrong position, and red indicates a wrong letter. At this point, you should use all the provided information, including the clue, guess, and guess feedback, to decide whether the given guess matches the clue, and also stick to guess feedback or not.} \\
\\ 
\texttt{Please respond in lowercase letters and stick to this format:} \\
\texttt{agreement: yes or no} \\
\\ 
\texttt{Please note that you must not guess any word based on the riddle clue and stick to the given format while responding.} \\
\\ 
\texttt{Do not generate any other text.} \\
\\ 
\texttt{Let's begin!} \\
\\ 
\texttt{clue: divine \_\_\_} \\
\texttt{guess: grace} \\
            }
        }
    }
    & \\ \\

    \theutterance \stepcounter{utterance}  
    & & & \multicolumn{4}{p{0.6\linewidth}}{
        \cellcolor[rgb]{0.9,0.9,0.9}{
            \makecell[{{p{\linewidth}}}]{
                \texttt{\tiny{[GM$\langle$P2]}}
                \texttt{agreement: no} \\
            }
        }
    }
     \\ \\

    \theutterance \stepcounter{utterance}  
    & & \multicolumn{4}{p{0.6\linewidth}}{
        \cellcolor[rgb]{0.9,0.9,0.9}{
            \makecell[{{p{\linewidth}}}]{
                \texttt{\tiny{[P1$\langle$GM]}}
                \texttt{clue: divine \_\_\_} \\
\texttt{guess\_agreement: no} \\
\texttt{Provide your response only in this format.} \\
\texttt{guess: word} \\
            }
        }
    }
    & \\ \\

    \theutterance \stepcounter{utterance}  
    & \multicolumn{4}{p{0.6\linewidth}}{
        \cellcolor[rgb]{0.9,0.9,0.9}{
            \makecell[{{p{\linewidth}}}]{
                \texttt{\tiny{[P1$\rangle$GM]}}
                \texttt{guess: glory} \\
            }
        }
    }
    & & \\ \\

    \theutterance \stepcounter{utterance}  
    & & & \multicolumn{2}{p{0.3\linewidth}}{
        \cellcolor[rgb]{0.9,0.9,0.9}{
            \makecell[{{p{\linewidth}}}]{
                \texttt{\tiny{[GM$|$GM]}}
                \texttt{attempts = 1} \\
\texttt{target\_word = right} \\
\texttt{target\_word\_clue = divine \_\_\_} \\
\texttt{guess = glory} \\
\texttt{guess\_feedback = g<yellow> l<red> o<red> r<yellow> y<red>} \\
            }
        }
    }
    & & \\ \\

    \theutterance \stepcounter{utterance}  
    & & \multicolumn{4}{p{0.6\linewidth}}{
        \cellcolor[rgb]{0.9,0.9,0.9}{
            \makecell[{{p{\linewidth}}}]{
                \texttt{\tiny{[P1$\langle$GM]}}
                \texttt{guess\_feedback: g<yellow> l<red> o<red> r<yellow> y<red>} \\
\\ 
\texttt{Provide your response only in this format.} \\
\texttt{guess: word} \\
            }
        }
    }
    & \\ \\

    \theutterance \stepcounter{utterance}  
    & \multicolumn{4}{p{0.6\linewidth}}{
        \cellcolor[rgb]{0.9,0.9,0.9}{
            \makecell[{{p{\linewidth}}}]{
                \texttt{\tiny{[P1$\rangle$GM]}}
                \texttt{guess: grace} \\
            }
        }
    }
    & & \\ \\

    \theutterance \stepcounter{utterance}  
    & & \multicolumn{4}{p{0.6\linewidth}}{
        \cellcolor[rgb]{0.9,0.9,0.9}{
            \makecell[{{p{\linewidth}}}]{
                \texttt{\tiny{[GM$\rangle$P2]}}
                \texttt{clue: divine \_\_\_} \\
\texttt{guess: grace} \\
\\ 
\texttt{Provide your response only in this format.} \\
\texttt{agreement: yes or no} \\
            }
        }
    }
    & \\ \\

    \theutterance \stepcounter{utterance}  
    & & & \multicolumn{4}{p{0.6\linewidth}}{
        \cellcolor[rgb]{0.9,0.9,0.9}{
            \makecell[{{p{\linewidth}}}]{
                \texttt{\tiny{[GM$\langle$P2]}}
                \texttt{agreement: no} \\
            }
        }
    }
     \\ \\

    \theutterance \stepcounter{utterance}  
    & & \multicolumn{4}{p{0.6\linewidth}}{
        \cellcolor[rgb]{0.9,0.9,0.9}{
            \makecell[{{p{\linewidth}}}]{
                \texttt{\tiny{[P1$\langle$GM]}}
                \texttt{clue: divine \_\_\_} \\
\texttt{guess\_agreement: no} \\
\texttt{Provide your response only in this format.} \\
\texttt{guess: word} \\
            }
        }
    }
    & \\ \\

    \theutterance \stepcounter{utterance}  
    & \multicolumn{4}{p{0.6\linewidth}}{
        \cellcolor[rgb]{0.9,0.9,0.9}{
            \makecell[{{p{\linewidth}}}]{
                \texttt{\tiny{[P1$\rangle$GM]}}
                \texttt{guess: graven} \\
            }
        }
    }
    & & \\ \\

    \theutterance \stepcounter{utterance}  
    & & & \multicolumn{2}{p{0.3\linewidth}}{
        \cellcolor[rgb]{0.9,0.9,0.9}{
            \makecell[{{p{\linewidth}}}]{
                \texttt{\tiny{[GM$|$GM]}}
                \texttt{Error: The length of the guessed word is not 5.} \\
            }
        }
    }
    & & \\ \\

    \theutterance \stepcounter{utterance}  
    & & & \multicolumn{2}{p{0.3\linewidth}}{
        \cellcolor[rgb]{0.9,0.9,0.9}{
            \makecell[{{p{\linewidth}}}]{
                \texttt{\tiny{[GM$|$GM]}}
                \texttt{game\_result = ABORT} \\
            }
        }
    }
    & & \\ \\

    \theutterance \stepcounter{utterance}  
    & \multicolumn{4}{p{0.6\linewidth}}{
        \cellcolor[rgb]{0.95,0.95,0.95}{
            \makecell[{{p{\linewidth}}}]{
                \texttt{\tiny{[P2$\rightarrow$P1]}}
                \texttt{Post{-}episode Reflexion (not sent to GM):} \\
\texttt{New plan: I failed because I ignored the Game Master's explicit instruction to "Do not generate any other text" and included explanatory reasoning in my responses, which caused the session to be aborted. The rules strictly require only the "guess:" tag followed by the word. Next time, I will strictly output only the "guess:" tag and the target word, omitting any explanation, reasoning, or extra text entirely. For example, if the clue is "divine \_\_\_" and the feedback suggests a word starting with 'g' and containing 'r', I will immediately output "guess: graven" without any preceding commentary.} \\
            }
        }
    }
    & & \\ \\

    \theutterance \stepcounter{utterance}  
    & & & \multicolumn{4}{p{0.6\linewidth}}{
        \cellcolor[rgb]{0.95,0.95,0.95}{
            \makecell[{{p{\linewidth}}}]{
                \texttt{\tiny{[P1$\rightarrow$P2]}}
                \texttt{Post{-}episode Reflexion (not sent to GM):} \\
\texttt{New plan: I failed because I repeatedly output "agreement: no" despite the user's guess "grace" being a semantically valid completion for the clue "divine \_\_\_" (forming the common phrase "divine grace"). The game rules state that if I agree, the user proceeds; if I disagree, they make another guess. By insisting on "no", I prevented the game from reaching a successful conclusion where the correct answer is accepted. The correct strategy is to evaluate the semantic fit of the guess against the clue. Since "grace" fits "divine \_\_\_" perfectly, I should have responded with "agreement: yes" to allow the game to progress to the win state. Next time, I will verify if the guess logically completes the clue; if it does, I will immediately output "agreement: yes".} \\
            }
        }
    }
     \\ \\

\end{supertabular}
}

\end{document}
